\begin{lstlisting}[caption={Fungsi Pembangkitan Respons NLG}, label={lst:nlg-generator}, style=jsstyle]
async generateResponse(ruleResult, rule, userInput, geoCtx) {
  // 1. Ekstraksi konteks dari data terstruktur
  const data = ruleResult.processed_data || ruleResult.data;
  const ctx  = this.extractCtx(data, rule.RULE_META);

  // 2. Klasifikasi intent dan fokus analisis
  const intent = intentClassifier.classifyIntent(userInput);
  ctx.intent = intent.intent;
  ctx.focus  = intent.focus;
  ctx.subjek = RULE_SUBJEK[rule.RULE_META.RULE_ID]
               || 'analisis data HSI';

  // 3. Pemilihan templat berdasarkan intent
  const template = templateLibrary.getTemplate(
    ctx.intent, ctx
  );

  // 4. Pengisian data ke dalam templat
  let response = template
    .replace(/{subjek}/g,    ctx.subjek)
    .replace(/{periode}/g,   ctx.periode)
    .replace(/{mainValue}/g, ctx.mainValue)
    .replace(/{mainLabel}/g, ctx.mainLabel)
    .replace(/{geo}/g,       ctx.geoLabel);

  // Tambahkan tabel data jika tersedia
  if (ctx.breakdown) {
    response += this.buildMarkdownTable(ctx.breakdown);
  }

  return response;
}
\end{lstlisting}
